% !TEX root =./ECE343_Practice_main.tex
%
% L13 practice set (Smith Charts IV): shorted and open stubs, and the
% single-stub shunt tuner.

\fullwidth{\textbf{Learning Objective 3.6: } I can use the Smith Chart to solve
transmission line problems.
\\[4pt]\normalfont\small Show your work. Solve each part on a printed Smith
chart, then check it with the stub formulas. Write a \textbf{Documentation}
line at the end (\textbf{None} if you did not collaborate). All lines and stubs
are lossless with $\myz{0} = 50\ \ohms$ and $v_p = 2\times10^{8}\ \text{m/s}$.}

\begin{questions}

%%%%%%%%%%%%%%%%%%%%%%%%%%%%%%%%%%%%%%%%%
\question \textbf{Stub reactance.}
\begin{parts}

	\begin{minipage}[t]{0.58\linewidth}
		\part Find the input impedance of a shorted stub $\lambda/8$ long, and the
		lumped element it acts like.
		\begin{solution}[1in]
		\eq{ z_{sc} &= j\tan\beta\ell = j\tan 45^\circ = j1 \\ \myz{sc} &= +j50\ \ohms }
		It acts like an inductor.
		\end{solution}
	\end{minipage}\hfill%
	\begin{minipage}[t]{0.37\linewidth}
		\raggedleft \vspace{12pt}
		$\myz{sc}$ = \ansbox[1.3in]{$+j50\ \ohms$}
	\end{minipage}

	\begin{minipage}[t]{0.58\linewidth}
		\part Find the input impedance of an open stub $\lambda/8$ long.
		\begin{solution}[1in]
		\eq{ z_{oc} &= -j\cot\beta\ell = -j\cot 45^\circ = -j1 \\ \myz{oc} &= -j50\ \ohms }
		It acts like a capacitor.
		\end{solution}
	\end{minipage}\hfill%
	\begin{minipage}[t]{0.37\linewidth}
		\raggedleft \vspace{12pt}
		$\myz{oc}$ = \ansbox[1.3in]{$-j50\ \ohms$}
	\end{minipage}

	\begin{minipage}[t]{0.58\linewidth}
		\part Find the shortest shorted stub and the shortest open stub that add a
		normalized susceptance of $+j1$ in shunt.
		\begin{solution}[1.3in]
		\eq{
			\text{short: } -\cot\beta\ell &= 1, & \beta\ell &= 135^\circ, & \ell &= 0.375\lambda \\
			\text{open: } \tan\beta\ell &= 1, & \beta\ell &= 45^\circ, & \ell &= 0.125\lambda
		}
		\end{solution}
	\end{minipage}\hfill%
	\begin{minipage}[t]{0.37\linewidth}
		\raggedleft \vspace{12pt}
		short = \ansbox[1.3in]{$0.375\lambda$}

		\vspace{4pt}
		open = \ansbox[1.3in]{$0.125\lambda$}
	\end{minipage}

	\begin{minipage}[t]{0.58\linewidth}
		\part Repeat part (c) for a susceptance of $-j1$.
		\begin{solution}[1.1in]
		\eq{
			\text{short: } \cot\beta\ell &= 1, & \ell &= 0.125\lambda \\
			\text{open: } \tan\beta\ell &= -1, & \ell &= 0.375\lambda
		}
		\end{solution}
	\end{minipage}\hfill%
	\begin{minipage}[t]{0.37\linewidth}
		\raggedleft \vspace{12pt}
		short = \ansbox[1.3in]{$0.125\lambda$}

		\vspace{4pt}
		open = \ansbox[1.3in]{$0.375\lambda$}
	\end{minipage}

\end{parts}

\newpage
%%%%%%%%%%%%%%%%%%%%%%%%%%%%%%%%%%%%%%%%%
\question \textbf{Single-stub tuner.} The load $\myz{L} = 100 + j50\ \ohms$ from
the Lesson 12 practice set must be matched at $500\ \mhz$ with a single shunt
stub. Lesson 12 found that $d = 0.1988\lambda$ puts the admittance at
$y = 1 + j1$.
\begin{parts}

	\begin{minipage}[t]{0.58\linewidth}
		\part Find the wavelength on the line and the distance $d$ in centimeters.
		\begin{solution}[1in]
		\eq{
			\lambda &= \dfrac{v_p}{f} = \dfrac{2\times10^{8}}{5\times10^{8}} = 0.4\ \text{m} \\
			d &= 0.1988(40\ \text{cm}) = 7.95\ \text{cm}
		}
		\end{solution}
	\end{minipage}\hfill%
	\begin{minipage}[t]{0.37\linewidth}
		\raggedleft \vspace{12pt}
		$d$ = \ansbox[1.3in]{$7.95\ \text{cm}$}
	\end{minipage}

	\begin{minipage}[t]{0.58\linewidth}
		\part The tuner is built in coaxial cable. Choose the stub type and find
		its length in centimeters.
		\begin{solution}[1.2in]
		The stub must add $-j1$. A coaxial stub should be shorted, and part 1(d)
		gives $0.125\lambda$.
		\eq{ \ell = 0.125(40\ \text{cm}) = 5.0\ \text{cm} }
		\end{solution}
	\end{minipage}\hfill%
	\begin{minipage}[t]{0.37\linewidth}
		\raggedleft \vspace{12pt}
		$\ell$ = \ansbox[1.3in]{$5.0\ \text{cm}$, shorted}
	\end{minipage}

	\begin{minipage}[t]{0.58\linewidth}
		\part Find the second solution: $d$, $y$ there, and the open-stub length.
		\begin{solution}[1.4in]
		Continuing clockwise from the upper crossing at $63.4^\circ$ to the lower one
		at $-63.4^\circ$ adds $126.9^\circ$, or $0.1762\lambda$.
		\eq{
			d &= 0.1988\lambda + 0.1762\lambda = 0.375\lambda, \qquad y = 1 - j1 \\
			\ell_{oc} &= 0.125\lambda = 5.0\ \text{cm}
		}
		\end{solution}
	\end{minipage}\hfill%
	\begin{minipage}[t]{0.37\linewidth}
		\raggedleft \vspace{12pt}
		$d$ = \ansbox[1.3in]{$0.375\lambda$}

		\vspace{4pt}
		$\ell_{oc}$ = \ansbox[1.3in]{$0.125\lambda$}
	\end{minipage}

	\part Explain in one or two sentences why a tuner is usually a shunt stub
	rather than a series stub.
	\begin{solution}[0.9in]
	A shunt stub connects at a T junction without cutting the main line. A series
	stub has to be inserted into one conductor of the line, which is hard to
	build in coax or microstrip.
	\end{solution}

\end{parts}

\vspace{1.5em}
\noindent\textbf{Documentation:} \rule[-2pt]{0.6\linewidth}{0.4pt}

\end{questions}
